\documentclass[11pt]{article}
\usepackage[a4paper,margin=25mm]{geometry}
\usepackage{amsmath,amssymb,amsthm,mathtools,hyperref,xurl}
\setlength{\emergencystretch}{2em}
\newtheorem{theorem}{Theorem}
\newtheorem{lemma}[theorem]{Lemma}
\newtheorem{corollary}[theorem]{Corollary}
\newtheorem{proposition}[theorem]{Proposition}
\newcommand{\Z}{\mathbb Z}\newcommand{\N}{\mathbb Z_{\ge0}}
\newcommand{\Q}{\mathbb Q}
\title{An exact three-Pell index certificate\\and nine fewer quadratic coordinates}
\author{\(\mathbb Z\mathbb Q\) Diophantine Project}
\date{23 November 2026}
\begin{document}
\maketitle
\begin{abstract}
We prove from elementary Pell identities the precise three-Pell block
occurring in the emended Jones system.  Its asserted meaning is the
value of a Pell sequence at an odd prescribed index.  The proof needs
no source-universality hypothesis.  It also gives an equivalent bridge
with one fewer squaring operation, a generalization to arbitrary
indices at least two, and a signed-witness argument which removes
two positivity wrappers.  A single quadratic square equation replaces
a third three-coordinate positivity wrapper.  Together these save
nine coordinates in the previously compiled quadratic core, with no
integer predicate omitted.  This settles local arithmetic obligations;
the full universality of the historical polynomial system remains a
separate issue.
\end{abstract}

\section{Elementary facts used in the proof}

For an integer \(a\ge2\), put \(D=a^2-1\) and define
\[
 X_j(a)+Y_j(a)\sqrt D=(a+\sqrt D)^j \qquad(j\in\Z).
\]
Then \(X_{-j}=X_j\), \(Y_{-j}=-Y_j\), and
\[
 X_0=1,\quad X_1=a,\qquad Y_0=0,\quad Y_1=1.
\]
Both sequences satisfy \(Z_{j+2}=2aZ_{j+1}-Z_j\).
For nonnegative indices, \(X_j\) is strictly increasing and \(Y_j\)
is strictly increasing; moreover \(Y_j\ge j\), with strict inequality
for \(j\ge2\).  These statements follow immediately from the recurrence
or multiplication by \(a+\sqrt D\).

\begin{lemma}[Classification and two congruences]\label{classification}
Every positive integer solution \(x,y\) of
\(x^2-(a^2-1)y^2=1\) has the form \((X_j(a),Y_j(a))\) for a unique
integer \(j\ge1\).  Also
\[
 Y_j(a)\equiv j\pmod2,
 \qquad Y_j(A)\equiv j\pmod M \quad\hbox{if } A\equiv1\pmod M.
\]
For every positive integer \(M\), there is a positive \(j\) with
\(M\mid Y_j(a)\).
\end{lemma}
\begin{proof}
For the classification, if \(y=1\) then \(x=a\).  If \(y>1\),
multiply \(x+y\sqrt D\) by \(a-\sqrt D\).  The new integer
coordinates are
\[
             x'=ax-Dy,\qquad y'=ay-x.
\]
The norm is still one.  Since
\(\sqrt D\,y<x<ay\) and \(\sqrt D>a-1\), one has
\(x'>0\) and \(0<y'<y\).  Descent gives a power of
\(a+\sqrt D\), and strict monotonicity gives uniqueness.
The two congruences follow from the recurrence, reduced modulo two
and modulo \(M\), respectively.  Finally multiplication by
\(a+\sqrt D\) is a unit operation in the finite ring
\((\Z/M\Z)[T]/(T^2-D)\).  Some positive power is one, and
the coefficient of \(T\) in that power is \(Y_j(a)\).
\end{proof}

\begin{lemma}[A square-divisibility step down]\label{stepdown}
Let \(p,q\ge1\), and put \(c=Y_p(a)\).
If \(c^2\mid Y_q(a)\), then \(c\mid q\).
In fact \(pc\mid q\).
\end{lemma}
\begin{proof}
Write \(q=kp+r\), where \(0\le r<p\), and put \(d=X_p(a)\).
Modulo \(c\), multiplication by
\((d+c\sqrt D)^k\) gives
\[
               Y_q(a)\equiv d^kY_r(a)\pmod c.
\]
The norm identity gives \(\gcd(c,d)=1\), and
\(0\le Y_r(a)<Y_p(a)=c\).  Thus \(r=0\), so \(q=kp\).
Now the binomial expansion of \((d+c\sqrt D)^k\) gives
\[
               Y_q(a)\equiv k d^{k-1}c\pmod {c^2}.
\]
The assumed divisibility implies \(c\mid k\), proving the stronger
conclusion.  The argument also covers \(c=1\).
\end{proof}

\begin{lemma}[The trace residue determines an index]\label{trace}
Let \(q>p\ge1\), put \(F=X_q(a)\), and let \(t\in\Z\).
Then
\[
      X_t(a)\equiv X_p(a)\pmod F
      \quad\Longleftrightarrow\quad
      t\equiv p\ \hbox{or }-p\pmod {4q}.
\]
\end{lemma}
\begin{proof}
The double-angle identities give
\[
 X_{2q}=2X_q^2-1\equiv-1\pmod F,
 \qquad Y_{2q}=2X_qY_q\equiv0\pmod F.
\]
Therefore \(X_{2q+j}\equiv-X_j\pmod F\) for every integer \(j\).
Together with \(X_{-j}=X_j\), this reduces any residue \(X_t\)
to \(\sigma X_r\), where \(0\le r\le q\) and
\(\sigma\in\{1,-1\}\).  More precisely, write
\(t=2qk+s\), \(0\le s<2q\).  If \(s\le q\), use
\(r=s,\sigma=(-1)^k\); otherwise use
\(r=2q-s,\sigma=(-1)^{k+1}\).

If \(\sigma=1\), strict monotonicity and
\(0<X_p<F\), \(0<X_r\le F\), imply \(r=p\).
If \(\sigma=-1\) and \(r=q\), the residue is zero and cannot
equal \(X_p\).  If \(r<q\), then, since \(q\ge2\),
\[
       0<X_r+X_p\le2X_{q-1}<X_q=F.
\]
The strict final inequality follows from
\(X_q=2aX_{q-1}-X_{q-2}\) and
\(X_{q-2}<X_{q-1}\).  So the negative sign is impossible.
The cases \(r=p,\sigma=1\) in the preceding explicit reduction
are exactly \(t\equiv\pm p\pmod {4q}\).  The reverse implication
follows from the same double-angle identities.
\end{proof}

\section{The exact odd-index certificate}

\begin{theorem}\label{odd}
Let \(a,B,c,d\) be positive integers such that
\(a\ge2\), \(B\ge3\) is odd, and \(c>B\).
For either of the choices
\begin{align}
 G_{\rm old}&=a+F^2(d^2-a),\label{oldG}\\
 G_{\rm new}&=a+F^2(c+1-a),\label{newG}
\end{align}
the following assertions are equivalent:
\begin{enumerate}
\item \(c=Y_B(a)\) and \(d=X_B(a)\).
\item There are positive integers \(i,F,j,o\) such that
\begin{align}
 d^2&=(a^2-1)c^2+1,\label{pell1}\\
 F^2&=(a^2-1)i^2c^4+1,\label{pell2}\\
 (d+oF)^2&=(G^2-1)(B+jc)^2+1.\label{pell3}
\end{align}
\end{enumerate}
Thus replacing \eqref{oldG} by \eqref{newG} preserves the
existence of positive integer witnesses, in both directions.
\end{theorem}
\begin{proof}
Assume the three equations.  By Lemma~\ref{classification}, write
\[
    c=Y_p(a),\quad d=X_p(a),\qquad
    ic^2=Y_q(a),\quad F=X_q(a)
\]
with \(p,q\ge1\).  Since \(c>B\ge3\), one has \(p\ge2\),
\(c\ge Y_2(a)=2a\), and \(p<c\).
Lemma~\ref{stepdown} gives \(c\mid q\), so \(q\ge c>p\).
Both choices of \(G\) satisfy \(G>a\), and
\[
 d^2\equiv F^2\equiv1\pmod c,
 \qquad G\equiv1\pmod c,
 \qquad G\equiv a\pmod F.                         \tag{*}
\]
For \(G_{\rm new}\), positivity uses \(c+1-a>0\).

Apply Lemma~\ref{classification} to \eqref{pell3}, obtaining
\[
                 B+jc=Y_t(G),\qquad d+oF=X_t(G)
\]
for a positive integer \(t\).  The first two congruences in (*)
and Lemma~\ref{classification} imply \(t\equiv B\pmod c\).
The recurrence for \(X_t\) and the last congruence in (*) imply
\[
                  X_t(a)\equiv X_t(G)\equiv d=X_p(a)\pmod F.
\]
Lemma~\ref{trace} gives \(t\equiv\pm p\pmod {4q}\), hence
\(B\equiv\pm p\pmod c\).  As \(0<B,p<c\), either \(B=p\)
or \(B=c-p\).  But \(c\equiv p\pmod2\), whereas \(B\) is odd.
The second alternative is impossible, proving \(p=B\).

Conversely, assume \(c=Y_B(a),d=X_B(a)\).
One may take the explicit index \(q=Bc\): expanding
\((d+c\sqrt D)^c\) shows \(c^2\mid Y_q(a)\).
Set
\[
                   i=Y_q(a)/c^2,\qquad F=X_q(a).
\]
All are positive integers and satisfy the first two equations.
For either choice of \(G\), properties (*) hold and \(G>a\).
Put
\[
             H=Y_B(G),\qquad I=X_B(G),\qquad
             j=(H-B)/c,\qquad o=(I-d)/F.
\]
The congruences (*) make \(j,o\) integers.  They are positive:
\(Y_B(G)>B\) for \(B\ge2\), and \(X_B(G)>X_B(a)\)
because \(G>a\).  For the latter inequality one can use
\(X_B(A)=\bigl((A+\sqrt{A^2-1})^B+
(A+\sqrt{A^2-1})^{-B}\bigr)/2\), which is strictly increasing
for real \(A>1\).  Finally the norm identity for \(I,H\)
is exactly \eqref{pell3}.
\end{proof}

\begin{corollary}[One less squaring operation]
In any polynomial construction using the old block with these
hypotheses, the substitution \(G_{\rm old}\mapsto G_{\rm new}\)
is valid at the level of existential positive integer solutions.
If an additional coordinate storing \(d^2\) was used only to form
\(G_{\rm old}\), that coordinate and its defining quadratic row
can be removed.  The first Pell equation itself already uses the
quadratic monomial \(d^2\) and does not require a separate register.
\end{corollary}

\section{A generalization without the odd-index restriction}

\begin{theorem}\label{arbitrary}
Let \(a\ge2\), \(B\ge2\), \(c>B\), and \(d>0\) be integers.
Then \(c=Y_B(a)\), \(d=X_B(a)\) if and only if positive integer
witnesses \(i,F,j,o\) satisfy
\[
 d^2=(a^2-1)c^2+1,
 \qquad F^2=4(a^2-1)i^2c^4+1,
\]
\[
 (d+oF)^2=(G^2-1)(B+2jc)^2+1,
 \qquad G=a+F^2(2c+1-a).
\]
\end{theorem}
\begin{proof}
Follow the preceding proof, with \(Y_q(a)=2ic^2\).
Now \(F^2\equiv1\pmod {2c}\), so \(G\equiv1\pmod {2c}\).
The last Pell index satisfies \(t\equiv B\pmod {2c}\).
Lemma~\ref{stepdown} still gives \(c\mid q\), and
Lemma~\ref{trace} gives \(t\equiv\pm p\pmod {4q}\), hence
\(B\equiv\pm p\pmod {2c}\).
Since \(0<B,p<c\), the negative alternative is impossible
without a parity assumption.  Conversely one may choose \(q=2Bc\),
since the preceding binomial expansion and the doubling identity give
\(2c^2\mid Y_q(a)\).  Take the last index to be \(B\),
as before.  All the resulting quotients are positive integers.
\end{proof}

\section{Signed witnesses in a rational quadratic system}

\begin{lemma}[A signed trace residue]\label{signedtrace}
Under the hypotheses of Lemma~\ref{trace},
\[
 X_t(a)\equiv X_p(a)\text{ or }-X_p(a)\pmod {X_q(a)}
 \quad\Longleftrightarrow\quad
 t\equiv p\text{ or }-p\pmod {2q}.
\]
\end{lemma}
\begin{proof}
The positive residue is Lemma~\ref{trace}.  The identity
\(X_{t+2q}\equiv-X_t\pmod {X_q}\) turns the negative residue
into the positive one.  The four resulting classes modulo \(4q\)
are exactly the two displayed classes modulo \(2q\).
\end{proof}

\begin{proposition}[No sign guard on the final norm root]
\label{signedblock}
Let \(a\ge2\), let \(B\ge3\) be odd, and let \(c,d\in\Z\)
satisfy \(|c|>B\) and \(d^2=(a^2-1)c^2+1\).
Let \(i\in\Z\setminus\{0\}\), \(j,o\in\Z\), and \(F\in\Q\)
satisfy the second and third norm equations
\eqref{pell2}--\eqref{pell3}, with \(G\) given by either
\eqref{oldG} or \eqref{newG}.  No sign assumption is made on
\(F\), \(B+jc\), or \(d+oF\).
Then
\[
                         |c|=Y_B(a),\qquad |d|=X_B(a).
\]
If \(c,d>0\), the projected pair is therefore exactly the same
as in Theorem~\ref{odd}, and positive witnesses can be reconstructed.
\end{proposition}
\begin{proof}
Set \(C=|c|\).  The first norm gives
\(C=Y_p(a), |d|=X_p(a)\) for some \(p\ge2\), so
\(C\ge2a\) and \(p<C\).  The second norm shows that the
rational number \(F\) has integer square.  Hence \(F\) is an
integer, and \(F\ne0\).  Classifying its absolute values gives
\[
                  |i|C^2=Y_q(a),\qquad |F|=X_q(a),
\]
where \(C\mid q\) and \(q\ge C>p\).

For the old bridge \(G>a\).  For the new bridge, if \(c>0\)
then \(G>a\); if \(c<0\), then
\[
 G=a-F^2(C+a-1)\le1-C\le-3.
\]
Thus \(|G|\ge2\) in every case.  Also
\(G\equiv1\pmod C\) and \(G\equiv a\pmod {|F|}\).
The quantity \(H=B+jc\) cannot be zero, because \(0<B<C\).
The last norm then also makes \(I=d+oF\) nonzero.
For the signed parameter \(G\), Pell classification says that
\[
          H=\varepsilon Y_t(G),\qquad I=\varepsilon'X_t(G)
\]
for some \(t\ge1\) and independent signs
\(\varepsilon,\varepsilon'\in\{1,-1\}\).
This follows from the classification at \(|G|\), using
\(X_t(-A)=(-1)^tX_t(A)\) and
\(Y_t(-A)=(-1)^{t-1}Y_t(A)\).

Reduction of the first relation modulo \(C\) gives
\(t\equiv\pm B\pmod C\).  Reduction of the second modulo
\(|F|\) gives \(X_t(a)\equiv\pm X_p(a)\pmod {|F|}\).
Lemma~\ref{signedtrace} yields \(t\equiv\pm p\pmod {2q}\),
so \(B\equiv\pm p\pmod C\).  The parity argument of
Theorem~\ref{odd} again excludes \(B=C-p\).  Thus \(p=B\).
\end{proof}

\section{The inequality makes the congruence and root positive}

\begin{proposition}[Automatic positivity of \(c-B\)]\label{phi}
Let \(r,s,w,n\) be positive integers, \(n\ge2\), and put
\[
 Y=sn^2,\qquad B=2r+1,\qquad
 a=(wn^2+1)rsn^2.
\]
Let \(p\) be either \(2ws^2r^2n^2\) or
\(2ws^2r^2n^6\).  Suppose that \(c,d,k,h\) are integers,
\(\tau\in\Q\), and
\begin{equation}\label{smallpell}
 \tau^2=(p^2-1)k^2+1,\qquad
 k=r+1+h(p-1),\qquad
 k^2-4(c-kY)^2>0.
\end{equation}
Suppose also that the signed three-Pell equations in
Proposition~\ref{signedblock} hold, without initially assuming
\(|c|>B\).  Then
\[
           k>0,\qquad c>B,\qquad h>0,
           \qquad c=Y_B(a),\quad |d|=X_B(a).
\]
In particular, if \(c=B+\varphi\), then \(\varphi>0\)
is automatic.
\end{proposition}
\begin{proof}
The defining formulas imply
\begin{equation}\label{bounds}
       Y\ge4,\qquad p\ge8r^2,\qquad 2a<p^2.
\end{equation}
For the last inequality it suffices to use the smaller value of \(p\).
After division by \(2rsn^2\), it reduces to
\[
        wn^2+1<2w^2s^3r^3n^2,
\]
which follows from \(wn^2\ge4\) and \(w,s,r\ge1\).
Also \(p-1>2(r+1)\).  The congruence for \(k\) therefore
implies \(|k|\ge r+1\).

The strict inequality in \eqref{smallpell} says
\[
                   \left|\frac ck-Y\right|<\frac12.
\]
Thus \(c\) and \(k\) have the same sign, and
\[
 |c|>|k|(Y-1/2)\ge(r+1)(Y-1/2)>2r+1=B.
\]
Proposition~\ref{signedblock} now applies and gives
\(|c|=Y_B(a)\), \(|d|=X_B(a)\).

Since \(\tau^2\) is an integer, \(\tau\) is an integer.
Classify the first norm in \eqref{smallpell} to obtain
\(|k|=Y_\ell(p)\) for some \(\ell\ge1\).
Suppose \(k<0\).  The congruence
\(Y_\ell(p)\equiv\ell\pmod {p-1}\) gives
\[
       \ell\equiv-(r+1)\pmod {p-1},
       \qquad \ell\ge p-r-2\ge4r+1.
\]
The final inequality follows from
\(p\ge8r^2\) and \(8r^2-5r-3\ge0\) for \(r\ge1\).
The elementary recurrence bounds give
\[
 |k|=Y_\ell(p)>p^{\ell-1}\ge p^{4r},
 \qquad |c|=Y_{2r+1}(a)<(2a)^{2r}<p^{4r}.
\]
This contradicts \(|c|>|k|(Y-1/2)>|k|\).  Hence \(k>0\),
and consequently \(c>0\), proving \(c>B\).

Finally \(k=r+1+h(p-1)>0\) and \(p-1>r+1\) force
\(h\ge0\).  If \(h=0\), then \(k=r+1\ge2\), \(k<2p\),
and
\[
         (pk-1)^2<(pk)^2-k^2+1<(pk)^2,
\]
contradicting that \(\tau^2\) is an integer square.
Thus \(h>0\) as well.
\end{proof}

\section{One square selects the positive first root}

\begin{lemma}[The odd-index square]\label{oddsquare}
For every integer \(r\ge0\),
\[
 X_{2r+1}(a)-1
       =(a-1)\bigl(Y_{r+1}(a)+Y_r(a)\bigr)^2.
\]
\end{lemma}
\begin{proof}
Write \(X=X_r(a),Y=Y_r(a)\).  Since
\(Y_{r+1}=X+aY\) and \(X^2=1+(a^2-1)Y^2\), expansion gives
\[
 (a-1)(Y_{r+1}+Y_r)^2
   =a-1+2(a^2-1)XY+2a(a^2-1)Y^2.
\]
The addition and doubling identities give the same expression for
\(X_{2r+1}-1\).
\end{proof}

\begin{proposition}[A one-coordinate replacement for \(\gamma>0\)]
\label{gamma}
Assume the hypotheses of Proposition~\ref{phi}, and also let
\(b\ge2\), \(n\ge b\), and \(\gamma\in\Z\) satisfy
\[
                 d=bw+(a-2)c+(4a-5)\gamma.
\]
Adjoin one rational coordinate \(\zeta\) and the quadratic equation
\begin{equation}\label{branch}
                         \zeta^2=(a-1)(d-1).
\end{equation}
Then \(d>0\) and \(\gamma>0\).
Conversely every positive solution of the old source block admits
a rational, in fact integral, \(\zeta\) satisfying \eqref{branch}.
\end{proposition}
\begin{proof}
Since \(a>1\), equation \eqref{branch} gives \(d\ge1\).
The norm and \(c>B\) exclude \(d=1\), so \(d>0\).
Proposition~\ref{phi} therefore gives
\(c=Y_B(a)\), \(d=X_B(a)\).
Put \(R_j=X_j(a)-(a-2)Y_j(a)\).  Multiplication gives
\[
 R_0=1,\quad R_1=2,\quad
 R_{j+1}=2R_j+(4a-5)Y_j(a),\quad R_2=4a-1.
\]
Thus \(R_B\ge R_2>a\).  Also
\[
                   a=(wn^2+1)rsn^2>wn^4\ge wb^4>wb.
\]
Consequently
\[
           \gamma=\frac{R_B-bw}{4a-5}>0.
\]
For completeness, Lemma~\ref{oddsquare} permits
\(\zeta=(a-1)(Y_{r+1}(a)+Y_r(a))\), proving the converse.
\end{proof}

\begin{corollary}[Nine coordinates removed from the compiled core]
Suppose a quadratic core has the displayed formulas and retains the
integer predicates for all its integer inputs, including the
nonzero integer \(i\), while ordinary positivity is implemented by
three-coordinate triangular-number wrappers.  The following changes
preserve its equivalence with the positive source system:
\begin{enumerate}
\item replace \(G_{\rm old}\) by \(G_{\rm new}\) and remove the
separate \(d^2\) register used only by the old bridge;
\item remove the positivity wrappers for \(I=d+oF\) and
\(\varphi=c-(2r+1)\);
\item replace the positivity wrapper for \(\gamma\) by
the single equation \eqref{branch} and its coordinate \(\zeta\).
\end{enumerate}
The net coordinate saving is \(1+3+3+(3-1)=9\).
The net equation saving is three: one square-register row and
two positivity rows disappear, while the third positivity row is
replaced by \eqref{branch}.  No integer predicate is discarded.
\end{corollary}
\begin{proof}
Every positive source solution admits the new branch coordinate by
Lemma~\ref{oddsquare}; the old and new bridge projections agree by
Theorem~\ref{odd}.  Conversely Proposition~\ref{phi} restores
\(c>B\) and \(h>0\), Proposition~\ref{gamma} restores \(d>0\)
and \(\gamma>0\), and Proposition~\ref{signedblock} identifies
the exact first Pell pair.  The constructive part of
Theorem~\ref{odd} then replaces \(i,F,j,o\) by positive
witnesses for the old block.  They occur outside the block only
in their integrality predicates and the defining arithmetic
registers, which can be chosen anew.  The remaining source
coordinates and their equations are unchanged.
\end{proof}

\section{Application to the complete rational quadratic tables}

The reduction has been applied to each complete stored system, including
the integral-input certificates and all their guards.  The resulting counts
are as follows; the first two rows begin with the global-contact-gauge
tables, and the last begins with the shared-reference ordinary-abc table.
\[
\begin{array}{c|rr|rr}
\text{integer interface}&\text{old variables}&\text{old rows}
 &\text{new variables}&\text{new rows}\\\hline
163&1185&1095&1176&1092\\
179&1355&1241&1346&1238\\
56\text{ (ordinary abc)}&846&728&837&725
\end{array}
\]
All residuals remain quadratic with integral parameter-polynomial
coefficients.  Their sum of squares is consequently a polynomial of total
degree four with the same rational solution set.  The instance parameters
occur in the coefficients and are not counted as existential unknowns.
The table compiler checks that the deleted square register occurs only in
its defining equation, the first norm, and the old bridge.  Each deleted
positivity coordinate occurs only in its own wrapper.  All other residuals
are preserved literally, up to coordinate numbering.

The complete systems are in
\nolinkurl{undecidability/coefficients/jones1980_pell_core/}.
They are produced by
\nolinkurl{undecidability/code/build_pell_core_reduction.py}
with option \texttt{--include-abc}.  The independent table checker
\nolinkurl{undecidability/verify_pell_core_reduction.py}
with the same option verifies every unchanged row, the three modified
formulas, coefficient domains, support, hashes and counts.  Its checks
are finite algebra; the existential equivalence follows from the proofs
above, together with the retained simultaneous integer-certificate
interfaces.

These counts are exact algebraic checkpoints.  Their interpretation as
undecidability bounds still assumes the explicitly unverified
source-universality assertion \(\mathcal U_{\rm em}\), as well as the
chosen integer interface's stated arithmetic premises.  The third row
also assumes ordinary rational abc.  No global minimum is claimed.

\section{The linear congruence and the remaining global obligation}

Put \(R_j=X_j(a)-(a-2)Y_j(a)\).  Direct multiplication gives
\[
 R_{j+1}-2R_j=(4a-5)Y_j(a),\qquad R_0=1.
\]
Thus the local certificate also gives the exact congruence
\[
         X_B(a)-(a-2)Y_B(a)\equiv2^B\pmod {4a-5}.
\]
Consequently, after the three-Pell block is established, the linear
row in the emended Jones system is precisely a positive-quotient
version of
\[
                      bw\equiv2^{2r+1}\pmod {4a-5}.
\]
Converting this congruence to an exact power relation requires the
remaining binomial-approximation inequalities and their size bounds.
The present theorem does not supply those bounds.

The earlier source audit also records an unresolved difference in
the auxiliary parameter \(p\).  The printed older display uses
\(p=2ws^2r^2n^2\) \cite[Theorem 3]{Jones80}.
Under the natural identification
\[
        U=2rsn^2,\qquad V=wn^2,\qquad A=2a=U(V+1),
\]
the modern BDHMS Pell/binomial bridge \cite[Definition 4.1 and
the square condition (4.2b)]{BDHMS} instead has parameter
\(U^2V/2=2ws^2r^2n^6\).  This is a concrete discrepancy in
the attempted identification, not a proof that the older system is
nonuniversal.  It must be resolved before importing a complete
universality assertion.  In particular, the elementary results here
do not discharge the source hypothesis \(\mathcal U_{\rm em}\),
and they do not justify replacing \(q=b^{5^{60}}\) by \(q=b^5\).

The accompanying exact checker
\nolinkurl{undecidability/verify_pell_index_optimization.py}
verifies the polynomial identities, produces explicit positive and
signed Pell witnesses, and checks a finite collection of trace residues.
Its finite examples are illustrations of the proved lemmas, not a
substitute for them.  It also verifies the separate negative-input
counterexample in \texttt{negative\_mixed\_witness.tex}.

\begin{thebibliography}{9}
\bibitem{Jones80}
J. P. Jones, \emph{Undecidable Diophantine equations},
Bulletin of the American Mathematical Society (New Series)
\textbf{3}, no.~2 (1980), 859--862, Theorem~3, p.~861.
\url{https://doi.org/10.1090/S0273-0979-1980-14832-6}.
\bibitem{BDHMS}
J. Bayer, M. David, M. Ha{\ss}ler, Y. Matiyasevich and D. Schleicher,
\emph{Diophantine equations over \(\Z\): universal bounds and parallel
formalization}, arXiv:2506.20909v3, 22 October 2026.
\url{https://arxiv.org/abs/2506.20909v3}.
\end{thebibliography}

\end{document}
